\chapter{Experimental Testbed and Evaluation Methodology}
\label{chap:empirical}

\section{Hardware and Containerized Testbed Configuration}
\label{sec:testbed}

To reflect a resource-constrained SME server node without external cloud network jitter, all experiments were executed on a dedicated local virtualization host configured with deterministic resource ceilings. Table~\ref{tab:env-specs} documents the exact testbed hardware, runtime, and instrumentation versions.

\begin{table}[h]
  \caption{Experimental testbed hardware, virtualization, and instrumentation specifications}
  \label{tab:env-specs}
  \centering
  \begin{tabular}{p{4.4cm}p{9.6cm}}
    \toprule
    \textbf{Component} & \textbf{Specification / Version} \\
    \midrule
    Host Processor & Intel Core i7-1165G7 (4 physical cores / 8 threads, up to 4.70 GHz) \\
    Host Operating System & Windows 11 (64-bit) with WSL2 Linux Kernel virtualization \\
    Allocated Virtual Memory & 5.0 GiB RAM ceiling, 4 virtual CPUs, 4.0 GiB swap (\texttt{.wslconfig}) \\
    Container Runtime & Docker Desktop / Docker Engine (Linux containers) \\
    Load Generator & Grafana \texttt{k6} CLI (HTTP/1.1 keep-alive virtual user simulation) \\
    Telemetry Collection & Docker \texttt{stats} daemon sampling, PowerShell stopwatch probes, Git diff analysis \\
    \bottomrule
  \end{tabular}
\end{table}

\section{Benchmarking Protocol and Workload Design}
\label{sec:benchmark-protocol}

The runtime evaluation protocol comprises four standardized measurement phases applied identically to both architectures:
\begin{enumerate}[leftmargin=1.5em]
  \item \textbf{Cold-to-Ready Initialization Time ($T_{\text{ready}}$):} Measures elapsed wall-clock seconds from container startup command invocation until the primary storefront/catalog HTTP endpoint returns status \texttt{200 OK}.
  \item \textbf{Idle Memory Footprint ($M_{\text{idle}}$):} Records steady-state container resident memory (MiB) after readiness stabilization with zero active virtual users ($0\text{ VUs}$).
  \item \textbf{Sustained Storefront Load Profile ($10\text{ VUs}$, $120\text{ s}$):} Executes a continuous 2-minute \texttt{k6} workload simulating concurrent shoppers browsing catalog listings and product details, capturing median latency ($p_{50}$), 90th-percentile latency ($p_{90}$), 95th-percentile tail latency ($p_{95}$), aggregate throughput (req/s), and HTTP failure rate.
  \item \textbf{Peak Concurrency Stress Memory ($M_{\text{load}}$, $20\text{ VUs}$):} Doubles concurrent virtual users to $20\text{ VUs}$ to measure peak container RAM allocation under saturation.
\end{enumerate}

\begin{algorithm}[h]
  \caption{Reproducible Benchmark Execution Protocol for Architecture $A$}
  \label{alg:benchmark}
  \begin{algorithmic}[1]
    \State Deploy stack $A \in \{\text{Monolith}, \text{Microservices}\}$ via Docker Compose
    \State Start high-resolution timer $t_0$; poll primary endpoint until HTTP $200$ at $t_1$
    \State Record $T_{\text{ready}} \gets t_1 - t_0$
    \State Wait $15\text{ s}$ for JIT and connection pool stabilization; sample $M_{\text{idle}}$ via \texttt{docker stats}
    \State Execute \texttt{k6 run --vus 10 --duration 120s} against catalog storefront route
    \State Extract $p_{50}, p_{90}, p_{95}$, total requests, req/s, and error percentage
    \State Execute \texttt{k6 run --vus 20 --duration 60s} and record peak resident memory $M_{\text{load}}$
    \State Tear down containers and purge transient network state before testing next stack
  \end{algorithmic}
\end{algorithm}

\section{Standardized Functional Evolution Scenarios}
\label{sec:change-scenarios}

To evaluate Hypothesis~H3 (maintainability and change propagation), five realistic SME e-commerce feature modifications were specified and traced across both repositories:
\begin{itemize}[leftmargin=1.5em]
  \item \textbf{Scenario S1 (Catalog Attribute Addition):} Add a structured \texttt{SustainabilityRating} field to catalog items and expose it in the product listing DTO.
  \item \textbf{Scenario S2 (Basket Promotional Discount Rule):} Implement a tiered percentage discount rule when basket subtotal exceeds a threshold.
  \item \textbf{Scenario S3 (Order Status Notification Hook):} Emit a domain/integration event when an order transitions to \texttt{Paid} status to trigger customer confirmation.
  \item \textbf{Scenario S4 (Payment Gateway Adapter Swap):} Replace the mock payment provider implementation with a new external gateway client interface.
  \item \textbf{Scenario S5 (Cross-Cutting Multi-Currency Display):} Introduce live exchange-rate conversion across catalog browsing, basket totals, and order summaries.
\end{itemize}

\section{Practitioner Interview Protocol}
\label{sec:interview-protocol}
To complement laboratory benchmarks with real-world organizational context, semi-structured interviews were conducted with four senior software engineers and technical leads (each with 6--12 years of production experience in .NET/cloud e-commerce systems). Questions focused on team size thresholds for service extraction, hidden operational failure modes, and debugging overhead.

\section{Chapter Summary}
\label{sec:ch3-summary}
This chapter detailed the 5.0~GiB WSL2/Docker testbed, formalized the four-stage runtime benchmarking algorithm, defined five standardized functional change scenarios (S1--S5), and outlined the practitioner interview protocol used to generate the empirical data in Chapter~4.
